We trained two proprioception-only policies with the same budget (1M decisions, one seed) on a waypoint task: a goal 2--3\,m away at a random bearing within $\pm60^\circ$ of the $+y$ axis, success when within 0.3\,m. The \emph{closed-loop} policy observes the goal position relative to the robot (2-d); the \emph{open-loop directional} policy observes only the unit direction to the goal at the start of the episode (no position feedback, no distance). Both are rewarded for speed toward the goal (the closed-loop one along the current bearing, the open-loop one along the initial bearing). We evaluate on 90 held-out terrains and goals each (amplitudes 0.04, 0.08, 0.12).

\textbf{Result (inconclusive; H3 not supported at this budget).} Neither policy reached the 0.3\,m goal radius within 8\,s (0/90 successes each), so time-to-goal cannot be compared. Using the reduction of goal distance instead, the open-loop directional policy was \emph{better}: it closed 51\% of the initial distance (mean final minimum distance 1.27\,m, 14\% of episodes within 0.6\,m) versus 30\% for the closed-loop policy (1.80\,m, 0\% within 0.6\,m); falls were 8\% and 2\%. This is the opposite of H3's prediction, but with speeds of roughly 0.1--0.2\,m/s a 2--3\,m goal is nearly out of reach in 8\,s, a single seed was trained, and the closed-loop reward changes direction as the robot moves, which slows credit assignment. We therefore report H3 as an unresolved stretch result and list the required follow-ups (longer episodes, goal-distance reward shaping, more seeds) in the next section.
